%! AP Statistics — Unit 8, Lesson 8.2 — Homework
\documentclass[10pt]{article}
\usepackage{apstats-article}
\usepackage{apstats-boxes}
\usepackage{grading}
\gradingassignment{apstats-u08-l02-hw}
\gradingcourse{advanced_statistics}

\begin{document}

\pageheader{Unit 8, Lesson 8.2}{Homework}
{\small\itshape Write each final answer in its box. Only what is inside a box is graded.}\par

\vspace{0.1in}
\textbf{Directions:} For each problem, set up a chi-square goodness of fit test as directed.
Show all work clearly.

\begin{enumerate}

  \item \textbf{Genetics.} A genetics researcher expects offspring to appear in a 9:3:3:1
        phenotype ratio. In a study, she observes 160 offspring with counts: Phenotype A = 92,
        Phenotype B = 28, Phenotype C = 26, Phenotype D = 14.

        \begin{enumerate}[label=\alph*.]
          \item State H$_0$ and H$_a$ for a chi-square goodness of fit test.

                H$_0$: \gradedblank[id=hw1a1, type=short, points=1, answer={p_A = 9/16, p_B = p_C = 3/16, p_D = 1/16}, check=text, display={$p_A = 9/16$, $p_B = p_C = 3/16$, $p_D = 1/16$ (9:3:3:1 ratio holds)}]{4.6in}{0.3in}

                H$_a$: \gradedblank[id=hw1a2, type=short, points=1, answer={at least one proportion differs from 9:3:3:1}, check=text, display={At least one proportion differs from 9:3:3:1.}]{4.6in}{0.3in}

          \item Calculate the expected count for each phenotype. (\textit{Hint:}
                Total ratio parts $= 9 + 3 + 3 + 1 = 16$.)

                \begin{center}
                \renewcommand{\arraystretch}{1.4}
                \begin{tabular}{lcccc}
                    \textbf{Phenotype} & A & B & C & D \\
                    \hline
                    Null proportion & $9/16$ & $3/16$ & $3/16$ & $1/16$ \\
                    \hline
                    Expected count & \makebox[2cm][r]{\smash{\raisebox{-3.5pt}{\gradedblank[id=hw1b1, type=short, points=1, answer=90, check=numeric, tolerance=0]{0.8cm}{0.2in}}}} & \makebox[2cm][r]{\smash{\raisebox{-3.5pt}{\gradedblank[id=hw1b2, type=short, points=1, answer=30, check=numeric, tolerance=0]{0.8cm}{0.2in}}}} & \makebox[2cm][r]{\smash{\raisebox{-3.5pt}{\gradedblank[id=hw1b3, type=short, points=1, answer=30, check=numeric, tolerance=0]{0.8cm}{0.2in}}}} & \makebox[2cm][r]{\smash{\raisebox{-3.5pt}{\gradedblank[id=hw1b4, type=short, points=1, answer=10, check=numeric, tolerance=0]{0.8cm}{0.2in}}}} \\
                \end{tabular}
                \end{center}

          \item Is the Large Counts condition met? Justify.
                \gradedhide{\writeline\\\gradedhere{\gradedblank[id=hw1c, type=rubric, points=2, display={Smallest expected 10 $\ge 5$: met.}]{\linewidth+6pt}{\gradedlinesheight{9}}}\writelines{9}}
        \end{enumerate}

  \boxguard[20]
  \item \textbf{Transportation Survey.} A local government claims that commuters use four
        modes of transportation in equal proportions (25\% each). In a random sample of
        200 commuters: car = 72, bus = 58, bike = 36, walk = 34.

        \begin{enumerate}[label=\alph*.]
          \item State H$_0$ and H$_a$ for a test of whether the proportions are equal.

                H$_0$: \gradedblank[id=hw2a1, type=short, points=1, answer={p_car = p_bus = p_bike = p_walk = 0.25}, check=text, display={$p_\text{car} = p_\text{bus} = p_\text{bike} = p_\text{walk} = 0.25$}]{4.6in}{0.3in}

                H$_a$: \gradedblank[id=hw2a2, type=short, points=1, answer={at least one mode's proportion is not 0.25}, check=text, display={At least one mode's proportion is not 0.25.}]{4.6in}{0.3in}

          \item Calculate the expected count for each transportation mode.

                $E = n \cdot p_0 = $ \smash{\gradedblank[id=hw2b, type=short, points=1, answer=50, check=numeric, tolerance=0]{6cm}{0.17in}}

                \begin{center}
                \renewcommand{\arraystretch}{1.3}
                \begin{tabular}{lcccc}
                    Mode & Car & Bus & Bike & Walk \\
                    \hline
                    Expected & \blank{1.8cm} & \blank{1.8cm} & \blank{1.8cm} & \blank{1.8cm} \\
                \end{tabular}
                \end{center}

          \boxguard[18]
          \item Check both conditions for a chi-square GOF test.
                \begin{itemize}
                    \item \textbf{Random:}
                          \gradedhide{\writeline\\\gradedhere{\gradedblank[id=hw2c1, type=rubric, points=1, display={Random sample stated.}]{\linewidth+6pt}{\gradedlinesheight{9}}}\writelines{9}}
                    \boxguard[15]
                    \item \textbf{Large Counts:}
                          \gradedhide{\writeline\\\gradedhere{\gradedblank[id=hw2c2, type=rubric, points=1, display={All expected $= 50 \ge 5$.}]{\linewidth+6pt}{\gradedlinesheight{9}}}\writelines{9}}
                \end{itemize}
        \end{enumerate}

  \item \textbf{True or False.} Write T or F in the T/F box and explain.

        ``If the sample size is large enough, we only need to check that each
        \emph{observed} count is greater than 5 for the chi-square GOF test.''

        \begin{center}
        \rule{0pt}{0.45in}\textbf{True} \qquad \textbf{False}\rlap{\hspace{0.5in}\textit{T/F:}~\smash{\gradedblank[id=hw3a, type=choice, points=1, answer=F]{0.4in}{0.2in}}}
        \end{center}

        Explanation: \gradedhide{\writeline\\\gradedhere{\gradedblank[id=hw3b, type=rubric, points=2, display={False: the condition is on \emph{expected} counts ($n p_0 \ge 5$), not observed counts.}]{\linewidth+6pt}{\gradedlinesheight{2}}}\writelines{9}}

\end{enumerate}

\end{document}
